\section{Métodos de recolección de datos para robots: real, teleoperación, simulación y síntesis}

La recolección de datos para robots puede seguir varias vías, y cada una tiene su costo y su sesgo. El despliegue real es el más cercano al producto, pero no alcanza la escala necesaria; la teleoperación permite obtener demostraciones de acciones, pero su costo en mano de obra es alto; la simulación permite escalar los escenarios de cola larga, pero tiene un sim-to-real gap; los videos de internet ofrecen gran escala, pero carecen de acciones y de retroalimentación. Un sistema de datos de alto nivel tiene que combinar estas fuentes.

\subsection{Comparación de los métodos de recolección}

\noindent\begin{longtable}{p{0.22\textwidth}p{0.30\textwidth}p{0.39\textwidth}}
\toprule
Método & Ventajas & Costos \\
\midrule
Despliegue real & Es el más cercano a las tareas y al hardware reales & Escala reducida, riesgo alto, pocos casos de cola larga. \\
Teleoperación & Demostraciones de acciones claras y alto control & Mano de obra costosa; la distribución de operadores es limitada. \\
Simulación & Puede generarse a escala y cubrir escenarios peligrosos o poco frecuentes & Brecha con la realidad y costo del modelado físico. \\
Datos sintéticos & Permiten completar con rapidez tareas y anotaciones de lenguaje & Sesgo de distribución y riesgo de alucinaciones. \\
Videos de internet & Escala enorme y cobertura de escenarios muy diversos & Carecen de acciones, de estados y de reward. \\
\bottomrule
\end{longtable}

\begin{knowledgebox}{Por qué los datos de fallos son los más valiosos}
Las trayectorias fallidas contienen condiciones límite e información de recuperación. Para un robot, las demostraciones exitosas le dicen al modelo «cómo hacerlo», mientras que las correcciones tras un fallo le dicen «dónde se equivoca» y «cómo recuperarse». Este tipo de datos tiene más valor marginal que repetir demostraciones exitosas.
\end{knowledgebox}

\subsection{Resumen de la sección}

La recolección de datos para robots no sigue una sola vía, sino que combina varias fuentes. El despliegue real, la teleoperación, la simulación, la síntesis y los videos de internet resuelven, cada uno, problemas distintos.
